\section{Experiments}\label{sec:experiments}
\paragraph{Data.}
\emph{HELM MMLU}, release v1.13.0 (summary date 2025-01-10): $79$ models on $57$ subjects, exact-match accuracy in points, no missing cells. \emph{Open LLM Leaderboard} (OLLB), the \texttt{contents} dataset at Hugging Face commit \texttt{9c09a7c}\ (2025-03-20; parquet file SHA-256 in Appendix~\ref{app:repro}): $\num{ollb.ceil.mean.ntar}+1$ models on six tasks (IFEval, BBH, MATH Lvl~5, GPQA, MuSR, MMLU-Pro) in the leaderboard's normalised points, after removing 79 duplicate names (the best variant is kept; the variant keeping the worst is a robustness slice) and one flagged model. The OLLB rule is the unweighted mean of the six normalised scores, which we verified reproduces its published average. We use only scores; neither source declares a licence for them. All numbers below are regenerated by the scripts in Appendix~\ref{app:repro}.

\paragraph{Pre-stated claims.}
Before running the grids the plan fixed six claims. Table~\ref{tab:claims} gives, for each, the refutation criterion as it was written, the measured value and the verdict. Where the plan fixed no number (C2, C3, C4, C6) any threshold we apply in the verdict was set after seeing the results (post hoc); C5 was reworded after the results, and marked as such.
\begin{table}[t]
\caption{Pre-stated claims: criterion fixed in the plan, measured value, verdict.}\label{tab:claims}
\centering
\fontsize{8}{9.5}\selectfont
\begin{tabular}{@{}p{0.3in} p{1.4in} p{2.3in} p{0.8in}@{}}
\toprule
 & Criterion in the plan & Measured & Verdict \\
\midrule
C1 & Closed forms (mean, minimum, median) match the solvers. & \num{theory.mismatch} mismatches in \num{theory.cells} cells (largest difference \num{theory.maxerr}). & Holds; a solver check. \\
C2 & $K$(mean)$\le K$(trimmed)$\le K$(median)$\le K$(min); refuted if the order fails ``for a sizeable share of targets'' (no number fixed). & $K$(median)$\le K$(min) at \textsf{ceil} for HELM \num{c2.helm.ceil.median<min}\%, OLLB \num{c2.ollb.ceil.median<min}\% of targets; OLLB at \textsf{gain10}, the three steps hold for \num{c2.ollb.gain10.mean<trimmed1}, \num{c2.ollb.gain10.trimmed1<median}, \num{c2.ollb.gain10.median<min}\%. & Refuted as a general order (post hoc: a step that fails for over 10\% of targets). \\
C3 & Almost every rule gameable at \textsf{ceil}; ``sharply'' different feasibility under caps; refuted if rule order and feasibility are stable across caps. & Feasible for 100\% of targets at \textsf{ceil}, \textsf{front0}, \textsf{front2}. At \textsf{gain10}: OLLB \num{ollb.gain10.trimmed1.feas}--\num{ollb.gain10.min.feas}\% by score rule, HELM \num{helm.gain10.min.feas}--\num{helm.gain10.plurality.feas}\% by rule. & Supported. \\
C4 & Rank-rule cost is not a monotone function of the score gap; ``reported either way''; no criterion. & At \textsf{ceil} Borda, the Condorcet requirement and plurality are cheaper in $L_1$ than the mean for \num{c4.helm.ceil.borda.cheaper}, \num{c4.helm.ceil.condorcet.cheaper} and \num{c4.helm.ceil.plurality.cheaper}\% of HELM targets; under caps they are infeasible more often (Section~\ref{sec:ordinal}). & Holds at \textsf{ceil}; qualified by feasibility. \\
C5 & Plan wording: cost for top targets is smaller than evaluation noise; first-place probability against investment to be reported. No threshold. & Post hoc restatement: a uniform raise by the nominal minimum $L_\infty$ puts the target first in a median of \num{r2.ollb.obs1.p0min}--\num{r2.ollb.obs1.p0max} (OLLB) and \num{r2.helm.obs1.p0min}--\num{r2.helm.obs1.p0max} (HELM) of noise draws over \num{r2.helm.obs1.ntar} targets per dataset; 90\% needs a further \num{r2.ollb.obs1.x90min_nomin}--\num{r2.ollb.obs1.x90max_nomin} points (all rules but minimum). & Plan wording not tested as written; restatement supported. \\
C6 & $K/n$ and $L_1$ against $n$, with and without cross-task correlation; no criterion. & Mean $L_1$ at $n=3$ and $n=57$: \num{nsweep.obs.mean.3.L1} and \num{nsweep.obs.mean.57.L1} observed, \num{nsweep.shuf.mean.57.L1} shuffled at $n=57$. & Descriptive; no claim that could fail. \\
\bottomrule
\end{tabular}
\end{table}

\subsection{Headline: cost by rule at the ceiling cap}\label{sec:headline}
Table~\ref{tab:summary} gives the median over targets of the three costs at the ceiling cap. Figure~\ref{fig:rank} shows the costs against the target's current mean rank under the \textsf{front0} cap, which is feasible for every target and binding, on the 136 OLLB targets on which all nine rules were run: $L_1$ rises with the original rank under every rule, fastest for the mean and slowest for the minimum, and $K$ stays between 1 and 6 tasks.
On OLLB the mean needs the \emph{fewest} tasks (median $K=\num{ollb.ceil.mean.K}$ of 6) and the \emph{most} points ($L_1=\num{ollb.ceil.mean.L1}$), against $L_1=\num{ollb.ceil.median.L1}$ for the median and $\num{ollb.ceil.min.L1}$ for the minimum; its $L_1$ exceeds the median's for all but \num{pw.ollb.L1.ceil.mean.median.cheaper}\% of targets and the minimum's for all but \num{pw.ollb.L1.ceil.mean.min.cheaper}\%.
On HELM (57 tasks) the mean needs $K=\num{helm.ceil.mean.K}$ and $L_1=\num{helm.ceil.mean.L1}$, the median $K=\num{helm.ceil.median.K}$ and $L_1=\num{helm.ceil.median.L1}$, the minimum $K=\num{helm.ceil.min.K}$ and $L_1=\num{helm.ceil.min.L1}$.
The cheapest rule depends on the measure: the minimum has the smallest $L_\infty$ on OLLB ($\num{ollb.ceil.min.Linf}$) but the largest on HELM ($\num{helm.ceil.min.Linf}$), where it must lift one weak subject far above the others.
The mean is dear in points because $L_1=n(\text{gap}+\varepsilon)$, the whole deficit summed over $n$ tasks, whereas the median and minimum repair only the tasks that matter to them (Appendix~\ref{app:theory}).
\paragraph{Is the gap between rules an artefact of score scale?} The rules' scores live on different scales, so we divide each rule's cost by its own leader gap or by the interquartile range of its score, and also match targets on rank under the rule (Appendix~\ref{app:newanalyses}). The ratio of the median $L_1$ to the mean's is on HELM \num{nr.helm.median.raw} (median rule) and \num{nr.helm.min.raw} (minimum) raw, and between \num{nr.helm.min.gap} and \num{nr.helm.median.gap} in all normalised and rank-matched views: five- to ten-fold. On the OLLB sample both rules are about three-fold cheaper (raw \num{nr.sample.median.raw}, \num{nr.sample.min.raw}), but the minimum's ratio rises to \num{nr.sample.min.iqr} by IQR. On the OLLB census only the minimum is several-fold cheaper (raw \num{nr.ollb.min.raw}, normalised \num{nr.ollb.min.gap}--\num{nr.ollb.min.iqr}); the median rule's ratio is \num{nr.ollb.median.raw} raw and \num{nr.ollb.median.gap}--\num{nr.ollb.median.iqr} normalised. The advantage is one of $L_1$ only: the ratio of median $L_\infty$ to the mean's is \num{nr.helm.median.rawinf} (median rule) and \num{nr.helm.min.rawinf} (minimum) on HELM and \num{nr.ollb.median.rawinf} and \num{nr.ollb.min.rawinf} on the census. Models with a zero on some task (\num{nozero.share}\% of OLLB) do not explain the minimum's low cost: removing them gives the ratio \num{nozero.ratio_removed} (\num{nozero.ratio_full} with them).

\begin{table}[t]
\caption{\capMainSummary}\label{tab:summary}
\centering
\fontsize{8}{9.5}\selectfont
\input{../tables/main_summary}
\end{table}

\begin{figure}[t]
\centering
\includegraphics[width=\textwidth]{cost_by_rank.pdf}
\caption{\capCostByRank}\label{fig:rank}
\end{figure}


\subsection{Caps decide feasibility}\label{sec:caps}
Figure~\ref{fig:capmap} is the cap map. At \textsf{ceil}, \textsf{front0} and \textsf{front2} every target is feasible for every rule, by construction or because the ceiling is 100. Under per-task gain caps the picture changes. With gains of at most 10 points per task, the share of OLLB targets that can reach first place is \num{ollb.gain10.mean.feas}\% under the mean and between \num{ollb.gain10.trimmed1.feas}\% (trimmed mean) and \num{ollb.gain10.min.feas}\% (minimum) over all score rules, and on HELM it ranges from \num{helm.gain10.min.feas}\% (minimum) to \num{helm.gain10.plurality.feas}\% (plurality) over all rules. With gains of at most 20 points the OLLB shares are \num{ollb.gain20.mean.feas}\% (mean) up to \num{ollb.gain20.min.feas}\% (minimum). With the empirical caps they are \num{ollb.emp95.mean.feas}\% (mean, \textsf{emp95}) and \num{ollb.emp99.mean.feas}\% (mean, \textsf{emp99}).
Costs at the ceiling are therefore no guide to realistic cost; under a cap a rule can look cheap only because it is feasible for the few targets already close, so medians are over feasible targets with the feasible share.
The cap also changes the cost ordering. On HELM the number of tasks of the mean is $\num{helm.ceil.mean.K}$ at \textsf{ceil} but $\num{helm.front0.mean.K}$ at \textsf{front0}, while its $L_1$ is unchanged ($\num{helm.ceil.mean.L1}$ and $\num{helm.front0.mean.L1}$), because the points needed are the same and only their spread over tasks is limited by the rivals' scores. Among the eight rules, the rank of the mean by median $K$ falls from \num{c3.helm.K.ceil.mean} (tied) at \textsf{ceil} to \num{c3.helm.K.front0.mean} at \textsf{front0}. C3 is supported: feasibility and ordering both depend on the cap.
\paragraph{A threshold instead of a grid.} The smallest uniform cap at which a target can reach first place is its $L_\infty$ at the ceiling (feasibility is monotone in the cap; this agrees with the stored 5-, 10- and 20-point solves in all \num{gs.valid.cells} (target, rule, cap) cells, \num{gs.valid.dis} disagreements). Call it $g^*$; the share of feasible targets at cap $g$ is the empirical distribution function of $g^*$ (Figure~\ref{fig:gstar}). The cap at which half of the targets can win ranges, over rules, from \num{gs.helm.crossmin} to \num{gs.helm.crossmax} points on HELM (mean \num{gs.helm.mean.cross}, 95\% interval \num{gs.helm.mean.lo}--\num{gs.helm.mean.hi}) and from \num{gs.ollb.crossmin} to \num{gs.ollb.crossmax} on the OLLB census (mean \num{gs.ollb.mean.cross}, \num{gs.ollb.mean.lo}--\num{gs.ollb.mean.hi}). The order of rules by this threshold differs between the leaderboards: the median rule is among the cheapest on HELM (\num{gs.helm.median.cross}) and the dearest on the OLLB census (\num{gs.ollb.median.cross}), the minimum the dearest on HELM (\num{gs.helm.min.cross}) and the cheapest on the census (\num{gs.ollb.min.cross}), and the order by $g^*$ differs from the order by $L_1$. For scale, among \num{gs.anchor.pairs} observed pairs of a fine-tune and its base model on the leaderboard, the 99th percentile of the gain is below the census median $g^*$ of the mean on \num{gs.anchor.q99.n} of the six tasks (at most \num{gs.anchor.q99.max5} points; IFEval \num{gs.anchor.q99.ifeval}); the pairs are a small selected subset, so this anchors the scale of $g^*$ and does not model what investment can buy.
\begin{figure}[t]
\centering
\includegraphics[width=\textwidth]{gstar_curves.pdf}
\caption{\capGstarCurves}\label{fig:gstar}
\end{figure}

\paragraph{Dependence between models.} OLLB models of one family are not independent targets. Resampling the \num{cb.nclusters} base-model families that contain the \num{ollb.ceil.mean.ntar} targets (\num{cb.singletons} singletons) widens the intervals by a median factor of \num{cb.width.med} (range \num{cb.width.min}--\num{cb.width.max}); the mean's median $L_1$ at the ceiling is \num{cb.mean.median_L1_ceil} with cluster interval \num{cb.mean.median_L1_ceil.clo}--\num{cb.mean.median_L1_ceil.chi} against \num{cb.mean.median_L1_ceil.nlo}--\num{cb.mean.median_L1_ceil.nhi} over targets. Dropping families of more than 10, 20 or 50 members changes no median $L_1$ by more than \num{cb.drop.maxrelL1}\% and no feasible share by more than \num{cb.drop.maxshare}. HELM has no family labels.
\begin{figure}[t]
\centering
\includegraphics[width=\textwidth]{cap_map.pdf}
\caption{\capCapMap}\label{fig:capmap}
\end{figure}


\subsection{Task-count orderings do not generalise}\label{sec:order}
C2 formalises a task-count intuition from our own plan, not from a published statement: the more of the score vector a rule uses, the fewer tasks it needs, $K(\text{mean})\le K(\text{trimmed})\le K(\text{median})\le K(\text{min})$. We use it as a test case of why task counts mislead. On OLLB at \textsf{ceil} the first two inequalities hold for \num{c2.ollb.ceil.mean<trimmed1}\% and \num{c2.ollb.ceil.trimmed1<median}\% of the \num{c2n.ollb.ceil.mean<trimmed1} targets, and the last for \num{c2.ollb.ceil.median<min}\%. At \textsf{front0} the three shares drop to \num{c2.ollb.front0.mean<trimmed1}\%, \num{c2.ollb.front0.trimmed1<median}\% and \num{c2.ollb.front0.median<min}\%, and on HELM at \textsf{ceil} the mean-median and median-minimum inequalities hold for \num{c2.helm.ceil.mean<median}\% and \num{c2.helm.ceil.median<min}\% of targets. We treat C2 as refuted as a general order. Each rule has its own rival threshold, so the number of tasks to raise is not ordered by how strict a rule seems; likewise ``the median needs $\lceil n/2\rceil+1$ tasks'' and ``the minimum needs all $n$'' are wrong in general, since tasks already above the winning level count (Appendix~\ref{app:theory}).

\subsection{Rank rules against the mean}\label{sec:ordinal}
Rank rules ignore magnitude, so their cost is not a function of the score gap (C4). At the ceiling cap they are cheaper in points than the mean for nearly all targets: on HELM the median ratio of $L_1$ (rank rule over mean) is $\num{c4.helm.ceil.borda.ratio}$ for Borda, $\num{c4.helm.ceil.condorcet.ratio}$ for the Condorcet requirement and $\num{c4.helm.ceil.plurality.ratio}$ for plurality, cheaper for \num{c4.helm.ceil.borda.cheaper}\%, \num{c4.helm.ceil.condorcet.cheaper}\% and \num{c4.helm.ceil.plurality.cheaper}\% of targets. On the OLLB sample (\num{c4.ollb_sample.ceil.borda.ntar} targets, rank-stratified) the ratios are $\num{c4.ollb_sample.ceil.borda.ratio}$, $\num{c4.ollb_sample.ceil.condorcet.ratio}$ and $\num{c4.ollb_sample.ceil.plurality.ratio}$.
The picture reverses in feasibility. Under 10-point gains on the OLLB sample, plurality is infeasible for \num{c4.ollb_sample.gain10.plurality.meanonly} targets for which the mean is feasible, and feasible for \num{c4.ollb_sample.gain10.plurality.ordonly} where the mean is not; Condorcet is infeasible for \num{c4.ollb_sample.gain10.condorcet.meanonly} such targets. Borda, in contrast, is feasible for \num{c4.ollb_sample.gain10.borda.ordonly} targets where the mean is not. Whether a rank rule is the cheaper route therefore depends on the cap and on the rule.
The OLLB sample is a deterministic grid of ranks, not a random sample, so the sampling step carries no seed. Redrawing the tail ranks with three random seeds (\num{ss.plurality.nlo}--\num{ss.plurality.nhi} targets each, ranks 2--21 always included) leaves the median $L_1$ ratio at \num{ss.borda.ratiolo}--\num{ss.borda.ratiohi} for Borda, \num{ss.plurality.ratiolo}--\num{ss.plurality.ratiohi} for plurality and \num{ss.condorcet.ratiolo}--\num{ss.condorcet.ratiohi} for the Condorcet requirement; the rank rule is cheaper for at least \num{ss.condorcet.cheaperlo}\% of targets; and at \textsf{gain10} the number of targets where plurality is infeasible but the mean feasible stays at \num{ss.plurality.ordlo}--\num{ss.plurality.ordhi} (Condorcet \num{ss.condorcet.ordlo}--\num{ss.condorcet.ordhi}); every redrawn interval overlaps the delivered one (Appendix~\ref{app:newanalyses}).
For the Condorcet requirement we report existence separately. On HELM a Condorcet winner exists on the observed matrix (it is the mean leader) and in \num{cond.helm.share}\% of \num{cond.helm.draws} noise draws; on OLLB it does not exist on the observed matrix nor in any of \num{cond.ollb.draws} draws, so there the Condorcet cost is the cost to beat every rival on a majority of tasks, not to unseat an existing winner (boundary cases in Table~\ref{tab:counter}, Appendix~\ref{app:boundary}).

\subsection{Near the top, evaluation noise matters}\label{sec:noise}
Near the top the investment is small: under the mean rule, the OLLB targets ranked 2 to 5 need between \num{near.ollb.L1.min} and \num{near.ollb.L1.max} points in total (HELM: \num{near.helm.L1.min} to \num{near.helm.L1.max}), and the scores carry sampling noise from finite test sets. We therefore ask not for the re-optimised cost under noise, which is a cost distribution, but how likely a \emph{fixed} investment is to work. We replay a uniform raise chosen on the observed leaderboard against $R=\num{r2.helm.obs1.R}$ noise draws of all scores, modelled as Gaussian with the binomial standard error of each cell, clipped to $[0,100]$ (Figure~\ref{fig:noise}b for one target; Figures~\ref{fig:replay-ollb} and \ref{fig:replay-helm} for six ranks between 1 and 20). The statistics below are medians over \num{r2.helm.obs1.ntar} targets per dataset (ranks 2 to 11 and 26 log-spaced ranks up to the last), with intervals over targets (Appendix~\ref{app:newanalyses}).
Raising every task by the nominal minimum $L_\infty$ puts the target first in a median of \num{r2.ollb.obs1.p0min} to \num{r2.ollb.obs1.p0max} of the draws on OLLB (score rules) and \num{r2.helm.obs1.p0min} to \num{r2.helm.obs1.p0max} on HELM, except that this probability is knife-edge, since the raise is designed to pass the best rival by $10^{-3}$ points. A more stable quantity is the further uniform raise needed to reach 90\%: a median of \num{r2.ollb.obs1.x90min_nomin} to \num{r2.ollb.obs1.x90max_nomin} points on OLLB and \num{r2.helm.obs1.x90min_nomin} to \num{r2.helm.obs1.x90max_nomin} on HELM for every rule but the minimum, and \num{r2.ollb.obs1.min.x90} (OLLB) and \num{r2.helm.obs1.min.x90} (HELM) for the minimum. Evaluation noise at the observed scale therefore adds a few points per task to the nominal minimum, and the nominal minimum is not a safe investment. A different noise seed or three redrawn target samples change the median probability by at most \num{r2.ollb.seed.dp0max} (OLLB) and \num{r2.helm.seed.dp0max} (HELM) and the extra raise by at most \num{r2.helm.seed.dx90max} points for the rules other than minimum and Borda.
The binomial standard error covers only item-sampling noise, and no repeated evaluations exist to estimate more, so 1.5 and 2 times the standard error are sensitivity multipliers, not estimates: the mean rule's extra raise for 90\% grows from \num{r2.helm.obs1.mean.x90} to \num{r2.helm.obs15.mean.x90} and \num{r2.helm.obs2.mean.x90} points on HELM and from \num{r2.ollb.obs1.mean.x90} to \num{r2.ollb.obs15.mean.x90} and \num{r2.ollb.obs2.mean.x90} on OLLB, roughly in proportion; the minimum rule's reaches \num{r2.ollb.obs2.min.x90} on OLLB and exceeds the 16-point grid for \num{r2.helm.obs2.min.cens} of \num{r2.helm.obs1.ntar} HELM targets, and only close pairs of other rules swap. Centring the noise on a shrunk truth (regression to the mean, a model with assumptions in Appendix~\ref{app:newanalyses}) \emph{raises} the probability (\num{r2.ollb.shr.p0min}--\num{r2.ollb.shr.p0max} on OLLB, \num{r2.helm.shr.p0min}--\num{r2.helm.shr.p0max} on HELM), so the observed-centred replay is the conservative reading. This supports C5 in its probabilistic form; cost below noise holds only for targets very near the top.
The sizes of all noise runs (draws, targets, seeds) are listed in Appendix~\ref{app:newanalyses}.
\begin{figure}[t]
\centering
\includegraphics[width=\textwidth]{scaling_and_noise.pdf}
\caption{\capScalingAndNoise}\label{fig:noise}
\end{figure}


\subsection{Number of tasks and cross-task correlation}\label{sec:nsweep}
On HELM, in 50 random subsets of $n\in\{3,\dots,57\}$ subjects, the mean's $L_1$ for a rank-5 target rises from $\num{nsweep.obs.mean.3.L1}$ at $n=3$ to $\num{nsweep.obs.mean.57.L1}$ at $n=57$, roughly in proportion to $n$, but only to $\num{nsweep.shuf.mean.57.L1}$ when each task column is permuted over models, which removes cross-task correlation; the median's $L_1$ falls with $n$ ($\num{nsweep.obs.median.12.L1}$ at $n=12$, $\num{nsweep.obs.median.57.L1}$ at $n=57$) and $K/n$ falls (Figure~\ref{fig:noise}a; details in Appendix~\ref{app:newanalyses}). This compares curves with and without correlation, not fitted exponents, and subsets of MMLU subjects are not independent leaderboards.

\subsection{Robustness}\label{sec:robust}
The rule order of the median $L_1$ at \textsf{ceil} (minimum $<$ median $<$ harmonic $<$ geometric $<$ mean on OLLB) is the same in all \num{slice.n_slices} data slices we ran (all models; merges removed; top 500 models; raw instead of normalised scores; duplicates resolved by keeping the worst variant), although the magnitudes change (for example the minimum's median $L_1$ is $\num{slice.ollb.min.L1}$ on all models and $\num{slice.ollb_raw.min.L1}$ on raw scores). A margin $\varepsilon$ up to 0.1 changes any median $L_1$ at \textsf{ceil} by at most \num{eps.maxrel.1e-1}\% (and the default $\varepsilon=10^{-3}$ by far less); the offset of the geometric and harmonic means (0.1 to 10 against 1) by at most \num{offs.maxrel}\%. Random Dirichlet weights, bootstrap over tasks and noise draws for the targets studied are in Appendix~\ref{app:robust}.
